% ══════════════════════════════════════════════════════════════════════
% Chapter 5: Misinterpretation and Cultural Distance
% ══════════════════════════════════════════════════════════════════════

\chapter{Misinterpretation and Cultural Distance}
\label{ch:misinterp}

\epigraph{``Understanding is always more than merely re-creating someone
else's meaning.''}{Hans-Georg Gadamer, \textit{Truth and Method}}

% ── Introduction ──────────────────────────────────────────────────────

The sender intends the receiver to compose signal~$s$ with
background~$r^*$ (the ``shared understanding'').  But the receiver
actually composes with~$r_{\mathrm{actual}}$---their real, possibly
different, cultural background.  The gap between
$r^* \compose s$ and $r_{\mathrm{actual}} \compose s$ is the formal
measure of \emph{misinterpretation}.

This chapter develops the formal theory of interpretive divergence.
Connecting to the Sapir--Whorf hypothesis~\cite{vygotsky1978}, which
holds that linguistic background shapes interpretation, and to Gadamer's
hermeneutic circle~\cite{gadamer1960}, which insists that understanding
requires pre-understanding, we show that cultural distance between
backgrounds governs interpretive distance between outcomes.

\medskip
\noindent\textbf{Chapter overview.}  We define the interpretation gap
(\S\ref{sec:ch5:gap}), analyse void backgrounds
(\S\ref{sec:ch5:void-bg}), develop iterated re-interpretation
(\S\ref{sec:ch5:iteration}), study resonance under iteration
(\S\ref{sec:ch5:resonance-iter}), and formalise cultural distance
(\S\ref{sec:ch5:distance}).

% ── §1. The Interpretation Gap ───────────────────────────────────────
\section{The Interpretation Gap}
\label{sec:ch5:gap}

\begin{definition}[Intended Interpretation]
\label{def:intended-interp}
The sender's \emph{intended interpretation} is
$\operatorname{intendedInterp}(r_{\mathrm{intended}}, s)
 = r_{\mathrm{intended}} \compose s$.
\end{definition}

\begin{definition}[Actual Interpretation]
\label{def:actual-interp}
The receiver's \emph{actual interpretation} is
$\operatorname{actualInterp}(r_{\mathrm{actual}}, s)
 = r_{\mathrm{actual}} \compose s$.
\end{definition}

\begin{definition}[Well-Interpreted]
\label{def:well-interpreted}
An interpretation is \emph{well-interpreted} if the actual and intended
interpretations resonate:
\[
  \operatorname{wellInterpreted}(r_{\mathrm{int}}, r_{\mathrm{act}}, s)
  \;\iff\;
  r_{\mathrm{int}} \compose s \;\res\; r_{\mathrm{act}} \compose s.
\]
This is weaker than equality but stronger than ``no relation.''
\end{definition}

\begin{theorem}[Resonant Backgrounds Guarantee Well-Interpretation]
\label{thm:resonant-well-interpreted}
If $r_{\mathrm{int}} \res r_{\mathrm{act}}$, then the signal~$s$ is
well-interpreted.
\end{theorem}

\begin{proof}[Proof sketch]
By resonance-compose-right applied to the resonant backgrounds.
See \texttt{resonant\_backgrounds\_well\_interpreted}.
\end{proof}

\begin{remark}
Members of the same culture (resonant backgrounds) always produce
well-interpreted outcomes for any signal.  Gadamer's ``fusion of
horizons'' is guaranteed when horizons resonate~\cite{gadamer1960}.
\end{remark}

\begin{theorem}[Self-Interpretation Always Succeeds]
\label{thm:self-interpretation}
$\operatorname{wellInterpreted}(r, r, s)$ for all $r, s$.
\end{theorem}

\begin{proof}[Proof sketch]
By reflexivity of resonance on the backgrounds.
See \texttt{self\_interpretation}.
\end{proof}

\begin{remark}
You always understand yourself.  Soliloquy is the degenerate case of
communication.
\end{remark}

\begin{theorem}[Well-Interpretation is Symmetric]
\label{thm:well-interp-symm}
If $\operatorname{wellInterpreted}(r_1, r_2, s)$, then
$\operatorname{wellInterpreted}(r_2, r_1, s)$.
\end{theorem}

\begin{proof}[Proof sketch]
By symmetry of resonance.  See \texttt{well\_interpretation\_symm}.
\end{proof}

\begin{remark}
Mutual intelligibility is symmetric.  If I can understand your ritual,
you can understand mine.  Cultural exchange is inherently
bidirectional---a formalisation of the anthropological observation that
first contact is always reciprocal.
\end{remark}

% ── §2. The Void Background ─────────────────────────────────────────
\section{The Void Background}
\label{sec:ch5:void-bg}

\begin{theorem}[Void Intended Transparency]
\label{thm:void-intended}
$\operatorname{intendedInterp}(\void, s) = s$.
\end{theorem}

\begin{proof}[Proof sketch]
By $\void \compose s = s$.  See \texttt{void\_intended\_transparency}.
\end{proof}

\begin{remark}
A sender who assumes zero shared context intends the signal to be
received ``as is.''  Any deviation from the raw signal is the receiver's
cultural contribution.  This is the formal model of the ``objective
journalist'' fantasy.
\end{remark}

\begin{theorem}[Void Actual Transparency]
\label{thm:void-actual}
$\operatorname{actualInterp}(\void, s) = s$.
\end{theorem}

\begin{proof}[Proof sketch]
By $\void \compose s = s$.  See \texttt{void\_actual\_transparency}.
\end{proof}

\begin{theorem}[Void-Void Perfect Interpretation]
\label{thm:void-void-perfect}
$\operatorname{wellInterpreted}(\void, \void, s)$ for all $s$.
\end{theorem}

\begin{proof}[Proof sketch]
Both interpret $s$ as $s$ itself, so the interpretations are identical
(and hence resonant).
See \texttt{void\_void\_perfect}.
\end{proof}

\begin{remark}
In the (impossible) limit of zero culture, there is zero
misinterpretation.  Misunderstanding is entirely a product of cultural
difference.
\end{remark}

\begin{theorem}[The Hermeneutic Circle]
\label{thm:hermeneutic-circle}
For any $r, s$:
$\operatorname{actualInterp}(r, s) \res \operatorname{actualInterp}(r, s)$.
\end{theorem}

\begin{proof}[Proof sketch]
By reflexivity of resonance.  See \texttt{hermeneutic\_circle}.
\end{proof}

\begin{remark}
Understanding always loops back to self-resonance.  You cannot escape
your own interpretive framework---even when you try to ``step outside,''
the result still resonates with your starting point.  This is the formal
content of Gadamer's hermeneutic circle~\cite{gadamer1960}.
\end{remark}

% ── §3. Iterated Re-Interpretation ───────────────────────────────────
\section{Iterated Re-Interpretation}
\label{sec:ch5:iteration}

\begin{definition}[Iterated Interpretation]
\label{def:iterate-interp}
The $n$-fold re-interpretation of signal~$s$ through background~$r$ is:
\[
  \operatorname{iterateInterp}(r, s, n) =
  \begin{cases}
    s & \text{if } n = 0, \\
    r \compose \operatorname{iterateInterp}(r, s, n-1) & \text{if } n > 0.
  \end{cases}
\]
This models rumination, echo chambers, and iterative cultural processing.
\end{definition}

\begin{lstlisting}[caption={Iterated interpretation in Lean~4},label=lst:iterate]
def iterateInterp (r s : I) : Nat -> I
  | 0 => s
  | n + 1 => IdeaticSpace.compose r (iterateInterp r s n)
\end{lstlisting}

\begin{theorem}[Zero Iteration = Raw Signal]
\label{thm:iterate-zero}
$\operatorname{iterateInterp}(r, s, 0) = s$.
\end{theorem}

\begin{proof}[Proof sketch]
By definition.  See \texttt{iterate\_zero}.
\end{proof}

\begin{theorem}[Single Iteration = Basic Interpretation]
\label{thm:iterate-one}
$\operatorname{iterateInterp}(r, s, 1) = r \compose s$.
\end{theorem}

\begin{proof}[Proof sketch]
By definition.  See \texttt{iterate\_one}.
\end{proof}

\begin{theorem}[Void Iteration Identity]
\label{thm:void-iterate}
$\operatorname{iterateInterp}(\void, s, n) = s$ for all $n$.
\end{theorem}

\begin{proof}[Proof sketch]
By induction on $n$, using $\void \compose x = x$ at the inductive step.
See \texttt{void\_iterate}.
\end{proof}

\begin{remark}
An ``empty mind'' processing a signal repeatedly never changes it.
Without cultural content, there is no interpretive drift.  Echo chambers
require content to echo.
\end{remark}

\begin{theorem}[Iterated Interpretation Depth Bound]
\label{thm:iterate-depth}
\[
  \depth(\operatorname{iterateInterp}(r, s, n))
  \leq n \cdot \depth(r) + \depth(s).
\]
\end{theorem}

\begin{proof}[Proof sketch]
By induction on $n$.  At step $n+1$, the depth-compose bound gives
$\depth(r \compose x) \leq \depth(r) + \depth(x)$; applying the
inductive hypothesis yields the result.
See \texttt{iterate\_depth\_bound}.
\end{proof}

\begin{remark}
Echo chambers have linearly bounded complexity.  Even infinite
rumination through a fixed cultural lens can only add linear depth.
Obsessive rumination is unproductive: each pass adds at most $\depth(r)$
more complexity.
\end{remark}

\begin{theorem}[Void Signal Single Iteration]
\label{thm:void-signal-iterate}
$\operatorname{iterateInterp}(r, \void, 1) = r$.
\end{theorem}

\begin{proof}[Proof sketch]
$r \compose \void = r$ by void-right.
See \texttt{void\_signal\_iterate\_one}.
\end{proof}

\begin{remark}
Processing ``nothing'' through a cultural lens yields the lens itself.
An empty news feed processed through conspiracy-theory schemas yields
the conspiracy schema unchanged.
\end{remark}

% ── §4. Resonance Under Iteration ────────────────────────────────────
\section{Resonance Under Iteration}
\label{sec:ch5:resonance-iter}

\begin{theorem}[Resonant Iteration Step]
\label{thm:resonant-step}
If $r_1 \res r_2$ and $x_1 \res x_2$, then
$r_1 \compose x_1 \;\res\; r_2 \compose x_2$.
\end{theorem}

\begin{proof}[Proof sketch]
By the full resonance-compose lemma.
See \texttt{resonant\_iteration\_step}.
\end{proof}

\begin{theorem}[Full Resonant Iteration]
\label{thm:full-resonant-iteration}
If $r_1 \res r_2$, then for all $n$:
\[
  \operatorname{iterateInterp}(r_1, s, n)
  \;\res\;
  \operatorname{iterateInterp}(r_2, s, n).
\]
\end{theorem}

\begin{proof}[Proof sketch]
By induction on $n$.  Base case: $s \res s$ by reflexivity.  Inductive
step: apply Theorem~\ref{thm:resonant-step} with the inductive
hypothesis.
See \texttt{resonant\_iteration}.
\end{proof}

\begin{remark}
Cultural similarity guarantees interpretive similarity at every stage
of processing.  This is why ``deep'' cultural understanding between
allied groups is possible---resonance propagates through all layers of
interpretation.
\end{remark}

\begin{theorem}[Self-Iteration Self-Resonance]
\label{thm:self-iter-resonance}
$\operatorname{iterateInterp}(r, s, n) \res
 \operatorname{iterateInterp}(r, s, n)$ for all $r, s, n$.
\end{theorem}

\begin{proof}[Proof sketch]
By reflexivity.  See \texttt{self\_iteration\_resonance}.
\end{proof}

\begin{remark}
The echo chamber's output is always ``consistent'' in the sense of
self-resonance---it just may not resonate with \emph{anything else}.
This is the formal trap of intellectual isolation.
\end{remark}

% ── §5. Cultural Distance ────────────────────────────────────────────
\section{Cultural Distance and Interpretation}
\label{sec:ch5:distance}

\begin{theorem}[Composition Preserves Well-Interpretation]
\label{thm:compound-well-interpreted}
If $s_1$ and $s_2$ are both well-interpreted under the same background
pair $(r_1, r_2)$, then
\[
  (r_1 \compose s_1) \compose (r_1 \compose s_2)
  \;\res\;
  (r_2 \compose s_1) \compose (r_2 \compose s_2).
\]
\end{theorem}

\begin{proof}[Proof sketch]
By the full resonance-compose lemma applied to the two well-interpretation
hypotheses.
See \texttt{compound\_well\_interpreted}.
\end{proof}

\begin{remark}
If two cultures understand each other's simple signals, they understand
compound signals too.  Mutual intelligibility at the atomic level
guarantees mutual intelligibility at the molecular level.  This is why
pidgins can bootstrap into creoles.
\end{remark}

\begin{theorem}[Interpretation Depth Gap]
\label{thm:interp-depth-gap}
\[
  \depth(\operatorname{actualInterp}(r_{\mathrm{actual}}, s))
  \leq \depth(r_{\mathrm{actual}}) + \depth(s).
\]
\end{theorem}

\begin{proof}[Proof sketch]
By depth-compose-le.  See \texttt{interpretation\_depth\_gap}.
\end{proof}

\begin{remark}
Misinterpretation does not create unbounded complexity.  Even the most
``wrong'' interpretation is bounded by the receiver's actual depth plus
the signal's depth.
\end{remark}

\begin{theorem}[Iteration Composition / Ritual Accretion]
\label{thm:iteration-composition}
\[
  \operatorname{iterateInterp}(r, s, n + m)
  = \operatorname{iterateInterp}(r,\;
    \operatorname{iterateInterp}(r, s, m),\; n).
\]
\end{theorem}

\begin{proof}[Proof sketch]
By induction on $n$.  See \texttt{iteration\_composition}.
\end{proof}

\begin{remark}
A lifetime of cultural processing ($n + m$ iterations) decomposes into
childhood processing ($m$~iterations) followed by adult processing
($n$~iterations).  Cultural development is \emph{sequentially
decomposable}: we can study ``what did early formation contribute?'' and
``what did later experience add?'' as separable phases.  This is the
formal basis of Bourdieu's distinction between primary and secondary
habitus formation~\cite{bourdieu1984}.
\end{remark}

% ── Summary ──────────────────────────────────────────────────────────
\section*{Chapter Summary}

This chapter developed the formal theory of misinterpretation within
IDT\@.  The \emph{interpretation gap} between intended and actual
interpretation is governed by the resonance between sender and receiver
backgrounds (Theorem~\ref{thm:resonant-well-interpreted}).
Well-interpretation is symmetric
(Theorem~\ref{thm:well-interp-symm}) and guaranteed for identical
backgrounds (Theorem~\ref{thm:self-interpretation}).  \emph{Iterated
re-interpretation} has linearly bounded depth
(Theorem~\ref{thm:iterate-depth}), and \emph{resonant iteration}
(Theorem~\ref{thm:full-resonant-iteration}) shows that culturally
similar backgrounds produce resonant interpretations at every stage.
The \emph{iteration composition} theorem
(Theorem~\ref{thm:iteration-composition}) provides the formal basis for
decomposing cultural development into sequential phases.

This concludes Part~I\@.  We have established the interpretive
framework: repertoires, signalling games, costly signals, resonance
activation, and misinterpretation.  Part~II develops the equilibrium
theory.

All eighteen theorems are verified in
\texttt{FormalAnthropology/IDT\_Signal\_Ch05.lean}.
